\documentclass[11pt]{article}
\usepackage[margin=1in]{geometry}
\usepackage{enumitem}
% =============================================================================
% Preambles/header.tex — shared LaTeX/Beamer preamble for this project
%
% Use from a Beamer lecture by prepending:
%     \input{header}
% and compiling with TEXINPUTS=../Preambles:$TEXINPUTS (the /compile-latex
% skill does this automatically).
%
% PALETTE CONTRACT. Color names defined here MUST match the names in
% Quarto/theme-template.scss so Beamer and Quarto renderings use the same
% palette. The scripts/check-palette-sync.sh script enforces this.
%
% To customize: edit the HEX values below AND the matching $variable in
% Quarto/theme-template.scss, then run scripts/check-palette-sync.sh.
% =============================================================================

% -----------------------------------------------------------------------------
% Engine detection + encoding
%
% Our /compile-latex skill uses XeLaTeX, where inputenc is unnecessary and
% can produce encoding warnings. Only load inputenc under pdfTeX.
% -----------------------------------------------------------------------------
\usepackage{iftex}
\ifPDFTeX
  \usepackage[utf8]{inputenc}
\fi

\usepackage{amsmath, amssymb, amsthm}
\usepackage{booktabs}
\usepackage{graphicx}

% -----------------------------------------------------------------------------
% PALETTE — mirrors Quarto/theme-template.scss variable names.
% Load xcolor and define colors BEFORE anything that references them
% (notably hyperref, hypersetup, and the Beamer color assignments).
% -----------------------------------------------------------------------------
\usepackage{xcolor}

\definecolor{primary-blue}{HTML}{012169}      % headings, accents
\definecolor{primary-gold}{HTML}{B9975B}      % emphasis, borders
\definecolor{highlight-yellow}{HTML}{F2A900}  % markers, alerts
\definecolor{light-bg}{HTML}{E8EDF5}          % title / section backgrounds
\definecolor{jet}{HTML}{1A1A1A}               % body text

% Semantic colors (used in diagrams and annotations)
\definecolor{positive}{HTML}{15803D}          % good / observed / identified
\definecolor{negative}{HTML}{B91C1C}          % bad / problematic / bias
\definecolor{neutral}{HTML}{525252}           % reference / context

% Highlight accents (match .hi-* classes in the SCSS)
\definecolor{hi-slate}{HTML}{314F4F}
\definecolor{hi-green}{HTML}{2E8B57}
\definecolor{hi-red}{HTML}{C41E3A}

% Semantic aliases so skill templates and snippets can write
% \textcolor{accent}{...} without hard-coding "primary-blue".
\colorlet{accent}{primary-blue}
\colorlet{accent2}{primary-gold}

% -----------------------------------------------------------------------------
% Hyperref — loaded AFTER xcolor + palette so link colors resolve.
% -----------------------------------------------------------------------------
\usepackage{hyperref}
\hypersetup{colorlinks=true, linkcolor=primary-blue, urlcolor=primary-blue,
            citecolor=primary-blue, pdfborder={0 0 0}}

% -----------------------------------------------------------------------------
% Course code — set once per deck (after \input{header}, before \begin{document})
% via \coursecode{CS401}. Shown in the footer on every slide so a deck's course
% is identifiable without opening the title page. Empty by default.
% -----------------------------------------------------------------------------
\makeatletter
\newcommand{\coursecode}[1]{\gdef\@coursecodetext{#1}}
\gdef\@coursecodetext{}
\makeatother

% -----------------------------------------------------------------------------
% Beamer theme assignments (only apply under Beamer).
%
% We detect Beamer via \@ifclassloaded{beamer}, which is the documented,
% non-internal way. This must be wrapped in \makeatletter/\makeatother
% because \@ifclassloaded contains an @-catcode character.
% -----------------------------------------------------------------------------
\makeatletter
\@ifclassloaded{beamer}{%
  \usefonttheme{professionalfonts}
  \setbeamercolor{structure}{fg=primary-blue}
  \setbeamercolor{frametitle}{fg=primary-blue}
  \setbeamercolor{title}{fg=primary-blue}
  \setbeamercolor{subtitle}{fg=primary-gold}
  \setbeamercolor{author}{fg=primary-blue}
  \setbeamercolor{institute}{fg=neutral}
  \setbeamercolor{date}{fg=primary-gold}
  \setbeamercolor{itemize item}{fg=primary-blue}
  \setbeamercolor{itemize subitem}{fg=primary-gold}
  \setbeamercolor{alerted text}{fg=primary-gold}
  \setbeamercolor{block title}{fg=white, bg=primary-blue}
  \setbeamercolor{block body}{bg=light-bg}
  % Semantic block variants: block=definition/concept (blue), exampleblock=
  % worked example (green/positive), alertblock=key takeaway or caution (gold).
  % Keep to <=2 colored boxes per slide across ALL three variants (INV-7).
  \setbeamercolor{block title example}{fg=white, bg=positive}
  \setbeamercolor{block body example}{bg=light-bg}
  \setbeamercolor{block title alerted}{fg=white, bg=primary-gold}
  \setbeamercolor{block body alerted}{bg=light-bg}

  % Minimal footer: course code (left) + page X / N (right), no navigation clutter
  \setbeamertemplate{navigation symbols}{}
  \setbeamertemplate{footline}{%
    \color{neutral}\footnotesize\hspace{0.7em}\@coursecodetext\hfill
    \insertframenumber/\inserttotalframenumber\hspace{0.7em}\vspace{0.3em}}
}{}
\makeatother

% -----------------------------------------------------------------------------
% TikZ setup — libraries the snippet gallery uses
% -----------------------------------------------------------------------------
\usepackage{tikz}
\usetikzlibrary{arrows.meta, positioning, calc, decorations.pathreplacing,
                fit, shapes.geometric, backgrounds}

% Shared TikZ styles — snippets and hand-written diagrams can pick these up
% via `\begin{tikzpicture}[every node/.style={...}]` or by naming specific
% styles (e.g., `\node[dag-node] (X) at (0,0) {X};`).
\tikzset{
  % Node styles for causal DAGs and similar
  dag-node/.style = {
    draw=primary-blue, fill=light-bg, rounded corners=2pt,
    minimum width=1.2cm, minimum height=0.9cm, align=center, font=\small
  },
  decision-node/.style = {
    draw=primary-gold, fill=white, diamond, aspect=1.8,
    minimum width=1.8cm, minimum height=1.2cm, align=center, font=\small,
    inner sep=1pt
  },
  % Edge styles — observed vs counterfactual
  observed-edge/.style = {-{Latex[length=2.5mm]}, thick, draw=jet},
  counterfactual-edge/.style = {-{Latex[length=2.5mm]}, thick, draw=neutral, dashed},
  confound-edge/.style = {-{Latex[length=2.5mm]}, thick, draw=negative, dashed},
  % Data-point styles for plots
  observed-dot/.style = {fill=primary-blue, draw=primary-blue, circle, inner sep=1.2pt},
  counterfactual-dot/.style = {fill=white, draw=neutral, circle, inner sep=1.2pt}
}

% -----------------------------------------------------------------------------
% Convenience macros
% -----------------------------------------------------------------------------
\newcommand{\muted}[1]{\textcolor{neutral}{#1}}
\newcommand{\key}[1]{\textcolor{primary-gold}{\textbf{#1}}}
\newcommand{\good}[1]{\textcolor{positive}{#1}}
\newcommand{\bad}[1]{\textcolor{negative}{#1}}

% Standout frame macro: full-bleed dark background for section transitions.
% Beamer-only (uses \begin{frame}/\setbeamercolor/\usebeamerfont) -- do not
% call from an article-class file (e.g. Notes/<CODE>/*-notes.tex). \muted,
% \key, \good, \bad, and \coursecode above are plain xcolor/gdef and are
% safe to use from either class; verified when Notes/article-class support
% was added.
% Expand: \transitionslide{Your transition title}
\newcommand{\transitionslide}[1]{%
  \begingroup
  \setbeamercolor{background canvas}{bg=primary-blue}
  \begin{frame}[plain]
    \centering
    \vfill
    {\usebeamerfont{title}\color{white}\Large #1\par}
    \vfill
  \end{frame}
  \endgroup
}

% Section divider frame (Beamer-only), an alternative to \transitionslide with a
% darker two-line style: near-black (jet) background, a small white label line
% above a larger gold title, and a thin gold rule along the bottom edge.
% Expand: \sectiondivider{Section 2}{Pointers}
\newcommand{\sectiondivider}[2]{%
  \begingroup
  \setbeamercolor{background canvas}{bg=jet}
  \begin{frame}[plain]
    \vspace*{3.0cm}
    {\color{white}\ttfamily\large #1\par}
    \vspace{0.5cm}
    {\color{primary-gold}\LARGE #2\par}
    \begin{tikzpicture}[remember picture, overlay]
      \draw[primary-gold, line width=2.5pt]
        ($(current page.south west)+(0,0.1)$) -- ($(current page.south east)+(0,0.1)$);
    \end{tikzpicture}
  \end{frame}
  \endgroup
}

% -----------------------------------------------------------------------------
% Channel watermark — "Concepts That Click" (YouTube).
%
% Article-class files (CompetitiveExam Q/A sets, Notes, Assignments) get a
% light diagonal draft watermark on every page plus a centered footer naming
% the channel. Beamer decks get a subtle TikZ background watermark instead
% (the footline already carries the course code). Centralized here so every
% MCQ PDF is branded without per-file edits.
% -----------------------------------------------------------------------------
\makeatletter
\@ifclassloaded{article}{%
  \usepackage{draftwatermark}
  \SetWatermarkText{Concepts That Click}
  \SetWatermarkScale{1}
  \SetWatermarkFontSize{2cm}
  \SetWatermarkLightness{0.86}
  \SetWatermarkAngle{45}
  \usepackage{fancyhdr}
  \pagestyle{fancy}
  \fancyhf{}
  \fancyfoot[C]{\footnotesize\textcolor{neutral}{Concepts That Click~|~YouTube: \href{https://www.youtube.com/@concepts.thatclick}{@concepts.thatclick}}}
  \renewcommand{\headrulewidth}{0pt}
  \renewcommand{\footrulewidth}{0pt}
  \fancypagestyle{plain}{%
    \fancyhf{}%
    \fancyfoot[C]{\footnotesize\textcolor{neutral}{Concepts That Click~|~YouTube: \href{https://www.youtube.com/@concepts.thatclick}{@concepts.thatclick}}}%
    \renewcommand{\headrulewidth}{0pt}%
    \renewcommand{\footrulewidth}{0pt}%
  }
}{}
\@ifclassloaded{beamer}{%
  \setbeamertemplate{background}{%
    \begin{tikzpicture}[remember picture, overlay]
      \node[rotate=45, opacity=0.055, align=center]
        at (current page.center)
        {\Huge\textcolor{neutral}{Concepts That Click}};
    \end{tikzpicture}%
  }
}{}
\makeatother 
\coursecode{JKSSB}
\renewcommand{\thesection}{05.\arabic{section}}
\newtheorem{example}{Example}[section]
\newenvironment{definitionbox}[1]{%
  \par\noindent\textbf{Definition (#1).}\ \itshape
}{\par\vspace{0.3em}}
\title{CPU, ALU, Control Unit \& Registers --- Detailed Lecture Notes}
\author{JKSSB Computer Awareness}
\date{\today}

\begin{document}
\maketitle

\noindent\textbf{Video lesson:}
\href{https://youtu.be/b6w8uT4GBYQ}{CPU Explained: ALU, Control Unit \& Registers}

\section{The CPU and its functional units}
The Central Processing Unit (CPU) is the part of the computer that executes
instructions. It repeatedly fetches an instruction from main memory, decodes
it, and then executes it. To do this work the CPU is organised into three
functional units: the Arithmetic Logic Unit (ALU), the Control Unit (CU),
and the registers. Data and control signals move between these units over
internal buses. The CPU works directly with main memory (RAM); it does not
reach into secondary storage such as a hard disk by itself.

\begin{definitionbox}{CPU (Central Processing Unit)}
The processor that fetches, decodes, and executes instructions, made up of
the ALU, the Control Unit, and registers.
\end{definitionbox}

A simple way to keep the roles apart is this: the ALU computes, the Control
Unit coordinates, and the registers hold the working values. Almost every
objective question about this topic turns on not swapping those three roles.

\section{The Arithmetic Logic Unit (ALU)}
The ALU is the computing part of the CPU. It performs two broad families of
operations. \textbf{Arithmetic} operations include addition, subtraction,
increment, decrement, multiplication, and division. \textbf{Logical}
operations include AND, OR, NOT, XOR, shifts, and comparisons. The ALU
operates on data held in binary form, and its operations are carried out with
logic-gate operations on those binary values. When the ALU finishes, it
returns its result to a register or the accumulator; it does not keep a
permanent copy.

\begin{definitionbox}{ALU (Arithmetic Logic Unit)}
The CPU unit that performs arithmetic and logical operations on binary data
using logic gates.
\end{definitionbox}

\begin{example}
For inputs $1$ and $1$, an addition produces the binary result for two, and
a logical XOR produces $0$ because the two bits are equal. Both results are
produced by the same unit --- the ALU --- and handed back for further use.
\end{example}

\section{The Control Unit (CU)}
The Control Unit directs and sequences the operations of the computer. It
decodes each instruction and then issues the control signals that time and
trigger the other units. It also controls the flow of data between the ALU,
the registers, and memory. What the Control Unit does \textbf{not} do is
compute: it does not add, subtract, or carry out logical operations. Those
are the ALU's job. Students frequently write ``the Control Unit performs the
calculations,'' which is exactly the error to avoid.

\begin{definitionbox}{Control Unit (CU)}
The CPU unit that directs and sequences operations, decodes instructions,
and issues control signals; it does not perform arithmetic or logic.
\end{definitionbox}

\begin{example}
When an addition instruction is executed, the Control Unit fetches and
decodes the instruction and signals the ALU to add the two values. The ALU
performs the addition; the Control Unit only directs the sequence.
\end{example}

\section{Registers}
Registers are the smallest and fastest storage in a computer. They are
located inside the CPU, so they are faster than cache and far faster than
main memory. Because they are very small, they hold only the values needed
for the work currently in progress: instructions, addresses, data, and
intermediate results. The following table lists the key registers and their
purposes.

\begin{center}
\begin{tabular}{p{0.30\textwidth}p{0.60\textwidth}}
\toprule
\textbf{Register} & \textbf{Purpose} \\
\midrule
Program Counter (PC / IP) & Holds the address of the \textbf{next} instruction to be fetched. \\
Instruction Register (IR) & Holds the instruction currently being decoded or executed. \\
Memory Address Register (MAR) & Holds the address of the memory location to be accessed. \\
Memory Data Register (MDR) & Holds the data being transferred to or from memory. \\
Accumulator (ACC) & Holds intermediate arithmetic and logic results. \\
General-purpose registers & Hold temporary operands and results. \\
Status / flags register & Holds condition flags such as zero, carry, sign, and overflow. \\
\bottomrule
\end{tabular}
\end{center}

The Program Counter is the register most often confused. It holds the
address of the \textbf{next} instruction to fetch, not the address of the
instruction currently being executed. The instruction currently being
decoded or executed is held in the Instruction Register. Memory access is
split between two registers: the MAR carries the address, and the MDR carries
the data itself. The accumulator is the natural home for an ALU result that
will be used again immediately.

\begin{example}
Suppose the CPU is about to fetch the instruction at address 200. The
Program Counter holds 200. Once the instruction is fetched, it is placed in
the Instruction Register for decoding. If the instruction needs to read a
value stored at address 900, the MAR holds 900 and the MDR holds the value
brought back from memory.
\end{example}

\section{Clock speed, cycles, cores, and threads}
The CPU's operations are paced by a clock. \textbf{Clock speed} is measured
in hertz (Hz); one gigahertz (GHz) is one billion clock cycles per second. A
\textbf{clock cycle} is a single tick of that clock. Do not assume every
instruction takes exactly one cycle: different instructions can require
different numbers of clock cycles. A \textbf{core} is an independent
processing unit within the CPU; a quad-core chip contains four cores. A
\textbf{thread} is a stream of instructions, and a system can run multiple
threads across the available cores.

\begin{definitionbox}{Clock cycle}
One tick of the system clock; the basic unit of timing for CPU operations.
\end{definitionbox}

\begin{example}
The office desktop in our running example is a quad-core machine running at
2.4 GHz. It has four independent cores, and each core can complete clock
cycles at about 2.4 billion per second. Running a spreadsheet calculation and
a browser at the same time lets different threads keep the cores busy.
\end{example}

\section{Exam retrieval and common confusions}
Three distinctions prevent most errors on this topic.
\begin{enumerate}
  \item Do not swap the ALU and the Control Unit. The ALU performs
    arithmetic and logic; the Control Unit only directs operations and
    issues control signals.
  \item Do not say the Program Counter holds the current instruction. It
    holds the address of the \textbf{next} instruction; the current
    instruction is in the Instruction Register.
  \item Do not place registers below cache or RAM in speed. Registers are
    the fastest storage, and the accumulator is where intermediate
    arithmetic and logic results are kept.
\end{enumerate}

\subsection*{Exercises}
\begin{enumerate}[label=\arabic*.]
  \item State the three main functional units of the CPU and one role of
    each.
  \item Expand ALU, CU, PC, IR, MAR, and MDR.
  \item Does the Control Unit perform arithmetic? Explain in one line.
  \item Which register holds the address of the next instruction, and which
    register holds the instruction currently being executed?
  \item Why are registers described as the fastest storage, and what does
    the accumulator hold?
\end{enumerate}

\subsection*{Solutions}
\begingroup\small
\begin{enumerate}[label=\arabic*.]
  \item ALU --- performs arithmetic and logical operations on binary data;
    Control Unit --- directs and sequences operations and issues control
    signals; registers --- hold instructions, addresses, data, and
    intermediate results.
  \item ALU = Arithmetic Logic Unit; CU = Control Unit; PC = Program
    Counter; IR = Instruction Register; MAR = Memory Address Register;
    MDR = Memory Data Register.
  \item No. The Control Unit directs and coordinates operations; the ALU
    performs the arithmetic and logical calculations.
  \item The Program Counter (PC/IP) holds the address of the next
    instruction; the Instruction Register (IR) holds the instruction
    currently being decoded or executed.
  \item Registers are located inside the CPU and are smaller and faster than
    cache and RAM. The accumulator (ACC) holds intermediate arithmetic and
    logic results.
\end{enumerate}
\endgroup
\end{document}
